\documentclass[12 pt, letterpaper]{hmcpset}
\usepackage{hw-preamble}

\name{Jayden Thadani}
\class{MATH 176 --- Algebraic Geometry}
\assignment{Homework assignment \#12}
\duedate{6th December 2024}

\begin{document}

For this week's set, let $X = \operatorname{Spec} \mathcal{O}$ for an integral domain $\mathcal{O}$. The \emph{quotient field} of $\mathcal{O}$ is the smallest field containing it, namely $K = \{ p/q \mid p, q \in \mathcal{O}, q \neq 0 \}$. For any Zariski open $U \subset X$, define the integral domain $\mathcal{O}_{X}(U) = \bigcap_{I \in U} \mathcal{O}_{I}$ in terms of the local rings $\mathcal{O}_{I} = \{ p / q \in K \mid q \notin I \}$.

\begin{problem}[Problem 1.]
	We say that subset $V \subset X$ is \emph{Zariski closed} if $V = Z_{J} = \{ I \in \operatorname{Spec} \mathcal{O} \mid J \subset I \} \cong \operatorname{Spec} \mathcal{O} / J$ for some ideal $J \subset \mathcal{O}$. We say that a subset $U \subset X$ is \emph{Zariski open} if $V = X \setminus U$ is Zariski closed.
	\begin{enumerate}
		\item Show that both $\O$ and $X$ are Zariski open.
		\item Say that $\{ U_{\alpha} \mid \alpha \in \mathcal{I} \}$ is an arbitrary family of Zariski open sets. Show that $U = \bigcup_{\alpha \in \mathcal{I}} U_{\alpha}$ is also Zariski open.
		\item Say that $\{ U_{\alpha} \mid \alpha \in \mathcal{I} \}$ is a \emph{finite} family of Zariski open sets. Show $U = \bigcap_{\alpha \in \mathcal{I}} U_{\alpha}$ is also Zariski open.
	\end{enumerate}
\end{problem}

\begin{solution}
	\begin{enumerate}
		\item Note that $\O$ is the set of primes $I \in \operatorname{Spec} \mathcal{O}$ that contain any maximal $M \in \operatorname{Spec} \mathcal{O}$ since, by definition, there are none. Thus, $\O$ is Zariski closed and its complement, $X$, is Zariski open.
			Similarly, note that $X = \operatorname{Spec} \mathcal{O} / \left < 0 \right >$, and thus, $X$ is Zariski closed too. Thus, its complement, $\O$, is Zariski open.
		\item To show that the arbitrary union of open sets is open, we need to it suffices to show their complement is closed. By De Morgan's laws, an arbitrary union of open sets $U_{\alpha}$ has complement
			\[
			X \setminus \left ( \bigcup_{\alpha \in \mathcal{I}} U_{\alpha} \right ) = \bigcap_{\alpha \in \mathcal{I}} X \setminus U_{\alpha}
			\]
			which is an intersection of closed sets.

			Write $V_{\alpha} = X \setminus U_{\alpha}$ for each $\alpha \in I$. By definition, $V_{\alpha} = \operatorname{Spec} \mathcal{O} / J_{\alpha}$ for some ideal $J_{\alpha} \subset \mathcal{O}$. Notice that if $J_{\alpha} \subset I \in \operatorname{Spec} \mathcal{O}$.
	\end{enumerate}
\end{solution}

\pagebreak

\begin{problem}[Problem 2.]
	Denote $\mathcal{O} = \mathbb{Z}$, so that its quotient field is $K = \mathbb{Q}$.
	\begin{enumerate}
		\item Assume that $I = \{ 0 \}$ is the zero ideal. Show that we have the localisation $\mathcal{O}_{I} = \mathbb{Q}$.
		\item Assume that $I = \left < \ell \right > = \ell \mathbb{Z}$ for a rational prime $\ell$. Show that $\mathcal{O}_{I} = \{ p / q \in \mathbb{Q} \mid \ell \text{ does not divide } q \}$.
	\end{enumerate}
\end{problem}

\begin{solution}
	\begin{enumerate}
		\item Since $\mathcal{O} = \mathbb{Z}$ is an integral domain, $I = \{ 0 \}$ is a prime ideal, and $\mathcal{O}_{I}$ is a well defined localisation away from $I$. In particular,
			\[
				\mathcal{O}_{I} = \{ p / q \mid q \notin I \} = \{ p / q \mid q \neq 0 \} = \mathbb{Q}.
			\]
		\item For primes $\ell \in \mathbb{Z}$, $I = \left < \ell \right >$ is a prime ideal of $\mathcal{O} = \mathbb{Z}$ consisting of the multiples of $\ell$. Thus, $\mathcal{O}_{I}$ is a well defined localisation away from $I$. In particular,
			\[
				\mathcal{O}_{I} = \{ p / q \mid q \notin I \} = \{ p / q \mid \ell \text{ does not divide } q \}.
			\]
	\end{enumerate}
\end{solution}

\pagebreak

\begin{problem}[Problem 3.]
	Since $\mathcal{O} = \mathbb{Z}$ is a Dedekind domain, we see that $X = \operatorname{Spec} \mathbb{Z}$ is a non-singular affine curve. Say that $V = \{ \ell_{1}, \ell_{2}, \dots, \ell_{n} \} \subset \mathbb{Z}$ is a finite collection of rational primes.
	\begin{enumerate}
		\item Upon identifying $\ell_{\alpha} \in V$ with $\left < \ell_{\alpha} \right > \in X$, show that the complement $U = X \setminus V$ is a Zariski open subset.
		\item Show that $\mathcal{O}_{X}(U) = \bigcap_{I \in U} \mathcal{O}_{I} = \{ p / q \in \mathbb{Q} \mid \text{no } \ell_{\alpha} \in V \text{ divides } q \}$.	
	\end{enumerate}
\end{problem}

\begin{solution}
	\begin{enumerate}
		\item Since each $\left < \ell_{\alpha} \right >$ is prime, each $V_{\alpha} = \{ \left < \ell_{\alpha} \right > \}$ is Zariski closed, and each $U_{\alpha} = X \setminus V_{\alpha}$ is Zariski open. By De Morgan's laws
			\[
				U = X \setminus \left ( \bigcup_{\alpha = 1}^{n} V_{\alpha} \right ) = \bigcap_{\alpha = 1}^{n} X \setminus V_{\alpha} = \bigcap_{\alpha = 1}^{n} U_{\alpha}.
			\]
			By the first exercise, this finite intersection of open sets is open.
		\item By the previous exercise, each localisation away from $\left < \ell_{\alpha} \right >$ is
			\[
				\mathcal{O}_{\left < \ell_{\alpha} \right >} = \{ p / q \in \mathbb{Q} \mid \ell_{\alpha} \text{ does not divide } q \}.
			\]
			Since $\mathcal{O}_{X}(U)$ is defined to be the intersection of all these local rings, we have
			\[
				\mathcal{O}_{X}(U) = \bigcap_{\alpha = 1}^{n} \mathcal{O}_{\left < \ell_{\alpha} \right >} = \{ p / q \mid \text{no } \ell_{\alpha} \in V \text{ divides } q \}.
			\]
	\end{enumerate}
\end{solution}

\pagebreak

\begin{problem}[Problem 4.]
	Denote $\mathcal{O} = F[x]$ as the collection of polynomials $p(x) = a_{n} x^{n} + \dots + a_{1} x + a_{0}$ with coefficients in a field $F$, so that its quotient field $K = F(x)$ consists of rational functions $p(x) / q(x)$.
	\begin{enumerate}
		\item Assume that $I = \{ 0 \}$ is the zero ideal. Show that the localisation $\mathcal{O}_{I} = F(x)$.
		\item Assume that $I = \left < f \right >$ for the linear polynomial $f(x) = x - r$. Show that $q(x) \in I$ if and only if $q(r) = 0$, and conclude that the localisation is $\mathcal{O}_{I} = \{ p(x) / q(x) \in F(x) \mid q(r) \neq 0 \}$.
	\end{enumerate}
\end{problem}

\begin{solution}
	\begin{enumerate}
		\item Since $I = \{ 0 \}$ is a prime ideal, we have the well-defined localisation
			\[
				\mathcal{O}_{I} = \{ p(x) / q(x) \mid q \notin I \} = \{ p(x) / q(x) \mid q(x) \neq 0 \} = F(x).
			\]
		\item Suppose $q(x) \in I$. Then $f$ divides $q$, so $q(r) = f(r) g(r) = 0$. Suppose conversely that $q(r) = 0$. Then by the division algorithm, $f(x) = x - r$  divides $q(x)$ and thus, $q \in I$.

			Since $I = \left < x - r \right >$ is a prime (in fact, it is maximal), we have the well defined localisation
			\[
				\mathcal{O}_{I} = \{ p(x) / q(x) \mid q \notin I \} = \{ p(x) / q(x) \in F(x) \mid q(r) \neq 0 \}
			\]
			where the last equality follows from our earlier result that $I$ contains exactly those polynomials vanishing at $r$.
	\end{enumerate}
\end{solution}

\pagebreak

\begin{problem}[Problem 5.]
	Since $\mathcal{O} = F[x]$ is a Dedekind domain, we see that $X = \operatorname{Spec} \mathcal{O} = \mathbb{A}^{1}_{F}$ is a non-singular affine curve. Say that $V = \{ r_{1}, r_{2}, \dots, r_{n} \} \subset F$ is a finite collection of elements.
	\begin{enumerate}
		\item Upon identifying $r_{\alpha} \in V$ with $\left < x - r_{\alpha} \right > \in X$, show the complement $U = X \setminus V$ is a Zariski open subset. Show that $V \cong Z_{J} = \{ I \in X \mid J \subset I \}$ in terms of the ideal $J = \left < (x - r_{1}) (x - r_{2}) \dotsm (x - r_{n}) \right >$.
		\item Show that $\mathcal{O}_{X}(U) = \bigcap_{I \in U} \mathcal{O}_{I} = \{ p(x) / q(x) \in F(x) \mid q(r_{\alpha}) \neq 0 \text{ for each } r_{\alpha} \in V \}$.
	\end{enumerate}
\end{problem}

\begin{solution}
	\begin{enumerate}
		\item Since each $\left < x - r_{\alpha} \right >$ is prime, each $V_{\alpha} = \{ \left < x - r_{\alpha} \right > \}$ is Zariski closed, and each $U_{\alpha} = X \setminus V_{\alpha}$ is Zariski open. By De Morgan's laws
			\[
				U = X \setminus \left ( \bigcup_{\alpha = 1}^{n} V_{\alpha} \right ) = \bigcap_{\alpha = 1}^{n} X \setminus V_{\alpha} = \bigcap_{\alpha = 1}^{n} U_{\alpha}.
			\]
			By the first exercise, this finite intersection of open sets is open.
		\item By the previous exercise, each localisation away from $\left < x - r_{\alpha} \right >$ is
			\[
				\mathcal{O}_{\left < x - r_{\alpha} \right >} = \{ p(x) / q(x) \in F(x) \mid q(r_{\alpha}) \neq 0 \}.
			\]
			Since $\mathcal{O}_{X}(U)$ is defined to be the intersection of all these local rings, we have
			\[
				\mathcal{O}_{X}(U) = \bigcap_{\alpha = 1}^{n} \mathcal{O}_{\left < x - r_{\alpha} \right >} = \{ p / q \mid q(r_{\alpha}) \neq 0 \text{ for each } r_{\alpha} \in V \}.
			\]
	\end{enumerate}
\end{solution}

\end{document}
